\section{掛金}

\subsection{掛金}

\begin{InyouJobunbox}
  \begin{itemize}
    \item 法第55条
    \item 令第35条
    \item 則第37条、第38条
    \item 解釈第4-1、第4-2、第8-13
  \end{itemize}
\end{InyouJobunbox}


\subsubsection{加入者掛金 〜法第55条第2項〜}

\begin{fillblank}[年金1 2021 問題1(1)A]
  \noindent \textbf{◯DB法}
  \begin{lawstructure}
    \begin{lawarticle}{第五十五条}{掛金}
      事業主は、給付に関する事業に要する費用に充てるため、規約で定めるところにより、年一回以上、定期的に掛金を拠出しなければならない。

      \lawparagraph{2}
      加入者は、政令で定める基準に従い規約で定めるところにより、前項の掛金の（　　A　　）することができる。

      （以下略）
    \end{lawarticle}
  \end{lawstructure}
\end{fillblank}

\begin{answerbox}
  \begin{enumerate}[A: ]
    \item 一部を負担
  \end{enumerate}
\end{answerbox}


\begin{truefalse}[年金1 2020 問題1(3)\textcircled{2}]
下線部の正誤を判定し、誤っている場合は正しい内容に改めなさい。
\\

  加入者が負担する掛金のある制度において、加入者がその選択により掛金を負担しない場合
の給付の額は、掛金を負担する加入者に比べて、当該\underline{掛金の拠出額の下限予定利率に基づく元
利合計額に相当}する程度の差を設けるものであること。
\end{truefalse}

\begin{answerbox}
  \textbf{誤}

  掛金の拠出額に相当
\end{answerbox}

\begin{truefalse}[年金1 2018 問題2(1)\textcircled{2}]
以下の文章の正誤を判定し、誤っている場合は正しい内容に改めなさい。
\\

掛金を負担していた加入者が、申し出により掛金を負担しないこととなった場合、規約の
変更によりその者が負担する掛金の額が減少することとなる場合を除き、当該加入者が掛金
を再び負担することはできない。
\end{truefalse}

\begin{answerbox}
  \textbf{正}
\end{answerbox}

\begin{shortanswer}[数理人試験 2025 問題4 設問3]
  確定給付企業年金法施行令第三十五条に定める加入者が掛金の一部を負担する場合の基準につい
て簡記せよ。
\end{shortanswer}

\begin{answerbox}
  \begin{enumerate}[(1)]
    \item 加入者が負担する掛金の額が当該加入者に係る掛金総額の二分の一を超えないこと
    \item 加入者が掛金を負担することについて、当該加入者の同意を得ること
    \item 掛金を負担している加入者が当該掛金を負担しないことを申し出た場合にあっては、当該
掛金を負担しないものとすること。
    \item 掛金を負担していた加入者であって、(2)もしくは(3)の規定により掛金を負担しないことと
なったものが当該掛金を再び負担することができるものでないこと（規約の変更によりその
者が負担する掛金の額が減少することとなる場合を除く）
  \end{enumerate}
\end{answerbox}

\begin{explanationbox}
  類題：年金１ 2024 問題2(1)ア
\end{explanationbox}

\begin{fillblank}[自作 2026/11/7]
    \noindent \textbf{◯DB令}
  \begin{lawstructure}
    \begin{lawarticle}{第三十五条}{加入者が掛金の一部を負担する場合の基準}
      法第五十五条第二項の政令で定める基準は、次のとおりとする。

      \lawitem{一}
加入者が負担する掛金の額が当該加入者に係る法第五十五条第一項に規定する掛金の額の（　　A　　）を超えないこと。
      
    \lawitem{二}
加入者が掛金を負担することについて、厚生労働省令で定めるところにより、当該加入者の同意を得ること。

      \lawitem{三}
掛金を負担している加入者が当該掛金を負担しないことを申し出た場合にあっては、当該掛金を負担しないものとすること。

      \lawitem{四}
掛金を負担していた加入者であって前二号のいずれかの規定により掛金を負担しないこととなったものが当該掛金を再び負担することができるものでないこと（規約の変更によりその者が負担する掛金の額が（　　B　　）することとなる場合を除く。）。

    \end{lawarticle}
  \end{lawstructure}

  \noindent \textbf{◯DB則}
  \begin{lawstructure}
    \begin{lawarticle}{第三十七条}{加入者が掛金を負担する場合の同意}
      令第三十五条第二号の加入者の同意は、規約で定めるところにより、加入者が掛金を負担することとなるとき及び規約の変更に伴い加入者が負担する掛金の額が（　　C　　）するときに得るものとする。
    \end{lawarticle}
  \end{lawstructure}


  \noindent \textbf{◯法令解釈}
  \begin{lawstructure}
    \begin{lawarticle}{}{}
      第４　掛金の額に関する事項

      1　加入者が負担する掛金に関する取扱い

      \lawitem{(1)}
      加入者が掛金を負担することについての同意は、規則第37条の規定により、加入者が掛金を負担することとなるとき及び規約の変更に伴い加入者が負担する掛金の額が増加するときに得ることとなっているが、この同意は、当該負担することとなる者又は負担額が増加する者について、（　　D　　）までに得なければならないこと。

      \lawitem{(2)}
      令第35条第4号の「掛金の額が（　　B　　）する」及び規則第37条の「掛金の額が（　　C　　）する」とは、規約の変更によって、変更後に加入者が負担することとなる掛金の額が変更前に加入者が負担していた掛金の額に比べて（　　B　　）又は（　　C　　）することをいい、掛金全体に占める加入者が負担する掛金の額の割合が（　　B　　）又は（　　C　　）する場合を指すものではないこと。

      \lawitem{(3)}
      加入者がその選択により掛金を負担しない場合においては、掛金を負担しない加入者の給付の額は、掛金を負担する加入者に比べて、（　　E　　）を設けるものであること。


    \end{lawarticle}
  \end{lawstructure}
\end{fillblank}

\begin{answerbox}
  \begin{enumerate}[A: ]
    \item 二分の一
    \item 減少
    \item 増加
    \item 掛金に係る規約の施行の日
    \item 掛金の拠出額に相当する程度の差
  \end{enumerate}
\end{answerbox}

\newpage

\subsubsection{掛金の額の算定方法 〜法第55条第3, 4項〜}

\begin{truefalse}[年金1 2023 問題1(3)イ]
  下線部の正誤を判定し、誤っている場合は正しい内容に改めなさい。
\\

  確定給付企業年金法施行規則第３８条第１項に定められている掛金の額の算定方法は、以下
のとおりである。

\begin{itemize}
  \item 加入者の給与に類するものに一定の割合を乗ずる方法
  \item \underline{加入者の年齢又は加入者が資格を取得したときの年齢}に応じて額を定める方法
  \item 加入者の給与又は給与に類するものに、（　　※　　）に応じて定めた割合を乗ずる方法
  \item 定額、給与に一定の割合を乗ずる方法及び上記の3 つの方法のうち２つ以上の方法を組み
合わせた方法
\end{itemize}

（注） （　　※　　）は出題の関係上省略している。
\end{truefalse}

\begin{answerbox}
  \textbf{誤}

  加入者の性別、年齢又は加入者が資格を取得したときの年齢
\end{answerbox}

\begin{fillblank}[数理人試験 2022 問題1 設問2]
  \noindent \textbf{◯DB則}
  \begin{lawstructure}
    \begin{lawarticle}{第三十八条}{掛金の額の算定方法}
      法第五十五条第四項第二号の厚生労働省令で定める適正かつ合理的な方法は、次のとおりとする。

      \lawitem{一}
      加入者の給与に類するものに（　　A　　）を乗ずる方法

      \lawitem{二}
      加入者の性別、（　　B　　）又は加入者が資格を取得したときの（　　B　　）に応じて額を定める方法

      \lawitem{三}
      加入者の給与又は給与に類するものに、加入者の性別、（　　B　　）又は加入者が資格を取得したときの
      （　　B　　）に応じて定めた割合を乗ずる方法

      \lawitem{四}
      定額、給与に（　　A　　）を乗ずる方法及び前三号の方法のうち（　　C　　）以上の方法を組み合わせた方法

      \lawparagraph{2}
      第四十五条第四項に規定するリスク分担型企業年金掛金額、第四十六条第一項に規定する特別掛金額、
      第四十六条の二第一項に規定するリスク対応掛金額、第四十七条の規定により計算される掛金の額、
      第五十二条第四項の規定により拠出する掛金の額及び第五十九条第一項の規定により掛金の額に追加して拠出する掛金の額は、
      前項の規定にかかわらず、それぞれ、第四十六条の三の規定により計算した額とする方法、
      第四十六条の規定により計算した額とする方法、第四十六条の二の規定により計算した額とする方法、
      第四十七条の規定により当該償却が次回の財政再計算のときに（　　D　　）するように計算された額とする方法、
      第五十二条第四項の規定により（　　E　　）の額から（　　F　　）を控除した額とする方法又は第五十九条第一項に
      規定する（　　G　　）とする方法により算定することができる。
    \end{lawarticle}
  \end{lawstructure}
\end{fillblank}

\begin{answerbox}
  \begin{enumerate}[A: ]
    \item 一定の割合
    \item 年齢
    \item 二
    \item 完了
    \item 数理債務
    \item 契約者価額
    \item 上回る額
  \end{enumerate}
\end{answerbox}

\begin{explanationbox}
  則第四十七条の掛金：次回財政計算までに発生する積立不足の予測額の償却のための掛金

  則第五十二条第四項の掛金：受託保証型DBを実施する場合に拠出できる特別掛金

  則第五十九条第一項の掛金：特例掛金
\end{explanationbox}

\newpage

\subsubsection{その他}


\begin{fillblank}[数理人試験 H26 問題1(3)a,b]
  \noindent \textbf{◯DB法}
  \begin{lawstructure}
    \begin{lawarticle}{第五十五条}{掛金}
      事業主は、（　　A　　）に要する費用に充てるため、規約で定めるところにより、年一回以上、定期的に掛金を拠出しなければならない。

      \lawparagraph{2}
      加入者は、政令で定める基準に従い規約で定めるところにより、前項の掛金の一部を負担することができる。

      \lawparagraph{3}
      掛金の額は、規約で定めるところにより算定した額とする。

      \lawparagraph{4}
      前項に規定する掛金の額は、次の要件を満たすものでなければならない。

      \lawitem{一}
      加入者のうち特定の者につき、不当に（　　B　　）なものであってはならないこと。

      \lawitem{二}
      定額又は給与に一定の割合を乗ずる方法その他適正かつ合理的な方法として厚生労働省令で定めるものにより算定されるものであること。
    \end{lawarticle}
  \end{lawstructure}
\end{fillblank}

\begin{answerbox}
    \begin{enumerate}[A: ]
    \item 給付に関する事業
    \item 差別的
  \end{enumerate}
\end{answerbox}

\newpage

\subsection{掛金の納付}



\begin{InyouJobunbox}
  \begin{itemize}
    \item 法第56条
    \item 令第36条
    \item 則第39条〜第42条
  \end{itemize}
\end{InyouJobunbox}


\begin{fillblank}[数理人試験 H26 問題1(3)c,d,e,f \& 年金1 2021 問題1(1)B]
  \noindent \textbf{◯DB法}
  \begin{lawstructure}
    \begin{lawarticle}{第五十六条}{掛金の納付}
      事業主は、前条第一項の掛金を、規約で定める日までに（　　A　　）に納付するものとする。

      \lawparagraph{2}
      事業主は、政令で定める基準に従い規約で定めるところにより、
      掛金を金銭に代えて金融商品取引法第二条第十六項に規定する金融商品取引所に上場されている
      （　　B　　）で納付することができる。ただし、事業主が当該（　　B　　）を（　　C　　）に納付する場合にあっては、
      当該（　　C　　）の同意を得たときに限る。
      
      \lawparagraph{3}
      （　　A　　）が、中小企業退職金共済法（昭和三十四年法律第百六十号）
      第十七条第一項又は第三十一条の四第一項の規定に基づき、
      （　　D　　）（第八十二条の五第一項及び第八十二条の六第一項において「機構」という。）
      から同法第十七条第一項に規定する厚生労働省令で定める金額の引渡し又は同法第三十一条の四第一項に
      規定する解約手当金に相当する額の移換を受けたときは、これらの金額については、
      前条及び第一項の規定により（　　E　　）とみなす。

    \end{lawarticle}
  \end{lawstructure}
\end{fillblank}


\begin{answerbox}
    \begin{enumerate}[A: ]
    \item 資産管理運用機関等
    \item 株式
    \item 基金
    \item 独立行政法人勤労者退職金共済機構
    \item 事業主が拠出した掛金
  \end{enumerate}
\end{answerbox}

\begin{explanationbox}
  資産管理運用機関等：資産管理運用機関または基金
\end{explanationbox}

\begin{fillblank}[自作 2026/11/7]
  \noindent \textbf{◯DB令}
  \begin{lawstructure}
    \begin{lawarticle}{第三十六条}{上場株式による掛金の納付}
      法第五十六条第二項の政令で定める基準は、次のとおりとする。

      \lawitem{一}
当該確定給付企業年金の（　　A　　）に当該確定給付企業年金に係る資産管理運用機関等が株式による掛金の納付を受けることができる旨の定めがあること。

\lawitem{二}
法第五十七条の掛金の額の基準に照らし（　　B　　）として厚生労働省令で定めるところにより算定される額の範囲内において行うものであること。

\lawitem{三}
納付する株式の価額は、（　　C　　）によるものとし、厚生労働省令で定めるところにより算定した額とすること。

\lawitem{四}
納付する株式の各銘柄につき、厚生労働省令で定めるところにより、前号の規定により算定した価額と当該確定給付企業年金に係る資産として既に運用されている株式（当該確定給付企業年金に係る資産以外の資産と合同して運用されているものを除く。次号において「既運用株式」という。）の価額との合計額が、当該確定給付企業年金に係る資産の総額の（　　D　　）に相当する額を超えないものであること。

\lawitem{五}
納付する株式の各銘柄につき、厚生労働省令で定めるところにより、納付する株式の数と当該確定給付企業年金に係る既運用株式の数との合計数が、発行済みの株式の総数の（　　E　　）を超えないものであること。

    \end{lawarticle}
  \end{lawstructure}
\end{fillblank}


\begin{answerbox}
    \begin{enumerate}[A: ]
    \item 規約
    \item 追加的に拠出すべき掛金
    \item 時価
    \item 百分の五
    \item 百分の五
  \end{enumerate}
\end{answerbox}

\newpage

\subsection{掛金の額の基準１}



\begin{InyouJobunbox}
  \begin{itemize}
    \item 法第57条
    \item 則第43条
    \item 解釈第4-3
    \item 死亡率告示
    \item 算定省令第3条
  \end{itemize}
\end{InyouJobunbox}


\begin{fillblank}[数理人試験 H26 問題1(3)g]
  \noindent \textbf{◯DB法}
  \begin{lawstructure}
    \begin{lawarticle}{第五十七条}{掛金の額の基準}
      掛金の額は、給付に要する費用の額の予想額及び予定（　　A　　）の額に照らし、
      厚生労働省令で定めるところにより、将来にわたって財政の均衡を保つことができるように計算されるものでなければならない。
    \end{lawarticle}
  \end{lawstructure}
\end{fillblank}

\begin{answerbox}
    \begin{enumerate}[A: ]
    \item 運用収入
  \end{enumerate}
\end{answerbox}

\begin{fillblank}[数理人試験 2025 問題1 設問2 A〜D]
  \noindent \textbf{◯DB則}
  \begin{lawstructure}
    \begin{lawarticle}{第四十三条}{掛金の額の計算に用いる基礎率及び財政悪化リスク相当額}
      法第五十七条に規定する掛金の額は、予定利率、予定死亡率、
      予定脱退率その他の通常予測給付額の算定の基礎となる率（以下「基礎率」という。）
      及び通常の予測を超えて（　　A　　）が損なわれる危険に対応する額として厚生労働大臣の定めるところにより算定した額
      （以下「財政悪化リスク相当額」という。）に基づき計算されるものとする。

      \lawparagraph{2}
      基礎率は、次のとおり定められるものとする。

      \lawitem{一}
予定利率は、（　　B　　）に基づき合理的に定められるものとする。
ただし、国債の利回りを勘案して厚生労働大臣が定める率を下回ってはならない。

\lawitem{二}
予定死亡率は、加入者等及びその遺族の性別及び年齢に応じた死亡率として
厚生労働大臣が定める率（以下「基準死亡率」という。）とする。
ただし、当該確定給付企業年金の加入者等及びその遺族の死亡の実績及び予測に基づき、
次の各号に掲げる加入者、（　　C　　）の区分に応じ、
当該各号に定める範囲内で定めた率を基準死亡率に乗じたものとすることができる。

\hangedlawitem{イ}{3}
加入者　零以上

\hangedlawitem{ロ}{3}
男子であって、（　　C　　）（ニに掲げる者を除く。）　〇・六八以上一・〇以下

\hangedlawitem{ハ}{3}
女子であって、（　　C　　）（ニに掲げる者を除く。　 〇・六五以上一・〇以下

\hangedlawitem{ニ}{3}
障害給付金の受給権者（イに掲げる者を除く。）　一・〇以上

      \lawitem{三}
予定脱退率は、当該確定給付企業年金の加入者の脱退の実績
（原則として、計算基準日の属する事業年度の前三事業年度の全部を含む三年以上の期間における実績とする。）
及び予測に基づき定められるものとする。

      \lawitem{四}
その他の基礎率は、当該確定給付企業年金における実績及び予測に基づき定められるものとする。

      \lawparagraph{3}
基礎率及び財政悪化リスク相当額は、財政計算ごとに定められるものとする。
ただし、前回の財政計算において定めた基礎率（（　　D　　）を除く。）
のうち継続して用いることが適切なものがある場合には、当該基礎率を継続して用いることができる。

    \end{lawarticle}
  \end{lawstructure}
\end{fillblank}

\begin{answerbox}
    \begin{enumerate}[A: ]
    \item 財政の安定
    \item 積立金の運用収益の長期の予測
    \item 加入者であった者又はその遺族
    \item 予定利率及び予定死亡率
  \end{enumerate}
\end{answerbox}


\begin{fillblank}[数理人試験 H25 問題1(3)a〜d]
  \noindent \textbf{◯DB則}
  \begin{lawstructure}
    \begin{lawarticle}{第四十三条}{掛金の額の計算に用いる基礎率及び財政悪化リスク相当額}
      法第五十七条に規定する掛金の額は、予定利率、予定死亡率、
      予定脱退率その他の通常予測給付額の算定の基礎となる率（以下「基礎率」という。）
      及び通常の予測を超えて財政の安定が損なわれる危険に対応する額として厚生労働大臣の定めるところにより算定した額
      （以下「財政悪化リスク相当額」という。）に基づき計算されるものとする。

      \lawparagraph{2}
      基礎率は、次のとおり定められるものとする。

      \lawitem{一}
      （略）

\lawitem{二}
予定死亡率は、加入者等及びその遺族の性別及び年齢に応じた死亡率として
厚生労働大臣が定める率（以下「基準死亡率」という。）とする。
ただし、当該確定給付企業年金の加入者等及びその遺族の死亡の実績及び予測に基づき、
次の各号に掲げる加入者、加入者であった者又はその遺族の区分に応じ、
当該各号に定める範囲内で定めた率を基準死亡率に乗じたものとすることができる。

\hangedlawitem{イ}{3}
加入者　（　　A　　）

\hangedlawitem{ロ}{3}
男子であって、加入者であった者又はその遺族（ニに掲げる者を除く。）　（　　B　　）

\hangedlawitem{ハ}{3}
女子であって、加入者であった者又はその遺族（ニに掲げる者を除く。　（　　C　　）

\hangedlawitem{ニ}{3}
障害給付金の受給権者（イに掲げる者を除く。）　（　　D　　）

（以下略）

    \end{lawarticle}
  \end{lawstructure}
\end{fillblank}

\begin{answerbox}
    \begin{enumerate}[A: ]
    \item 零以上
    \item 〇・六八以上一・〇以下
    \item 〇・六五以上一・〇以下
    \item 一・〇以上
  \end{enumerate}
\end{answerbox}



\begin{explanationbox}
  法改正に注意。
\end{explanationbox}


\begin{shortanswer}[年金2 2022 問題2(4)]
  「確定給付企業年金に関する数理実務ガイダンス」には、予定利率を設定する際に検討
すべき事項として、事業主等へ助言を行うことが望ましいとされている点が記載されてい
る。これを3 つ簡潔に入力しなさい。
\end{shortanswer}

\begin{answerbox}
  次のうちから3つを記載
  \begin{itemize}
    \item 長期的期待収益率および下限予定利率との整合性
    \item 設定した予定利率の年金財政への影響
    \item 適用後の決算時における利差損益の動向
    \item 長期的期待収益率とリスクとの整合性
    \item その他、年金財政への影響が大きいと思われる事項
  \end{itemize}
\end{answerbox}

\begin{shortanswer}[自作 2026/11/7]
  「確定給付企業年金に関する数理実務ガイダンス」には、
  予定利率の影響について事業主に助言を行うことが望ましいケースが２点例示されている。これを簡記しなさい。
\end{shortanswer}

\begin{answerbox}
  \begin{itemize}
    \item 資産構成のリスクが異常に大きく、財政運営上の支障が予想される場合
    
    \item 予定利率が直近の実勢利回り(直近の長期国債利回り等)を大幅に上回っている場合
  \end{itemize}
\end{answerbox}

\begin{truefalse}[年金2 H29 問題1(1)\textcircled{3}]
  下線部の正誤を判定し、誤っている場合は正しい内容に改めなさい。
\\

  開放基金方式・総合保険料方式（閉鎖型）・加入年齢方式は、\underline{発生給付評価方式}と呼ばれ、将
来の給付額を予測し、将来の掛金収入と年金資産の合計額が将来の給付額と一致するように
掛金を算定する。
\end{truefalse}

\begin{answerbox}
  \textbf{誤}

  予測給付評価方式
\end{answerbox}

\begin{shortanswer}[自作 2026/11/7]
  発生給付評価方式に分類される財政方式を一つ挙げなさい。
\end{shortanswer}

\begin{answerbox}
  予測単位積増方式
\end{answerbox}


\begin{shortanswer}[年金2 2025 問題2(5)]
  確定給付企業年金法施行規則第 43 条では、掛金の額の計算に用いる基礎率である予定利率を 次のとおり定めている。
\\

  予定利率は、積立金の（　　　　）に基づき合理的に定められるものとする。ただし、\underline{国債の利回りを勘案して}\\
  \underline{厚生労働大臣が定める率}を下回ってはならない。
\\

これについて以下の(1)、(2)、(3)に答えなさい。

\begin{enumerate}[(1)]
  \item 空欄にあてはまる語句を入力しなさい。(20 字以内)\\
  \item 下線部は、「確定給付企業年金制度について(平成 14 年 3 月 29 日年発第 0329008 号)」
  に下限予定利率としてその設定の考え方が示されている。この下限予定利率の設定の考え方を簡潔に入力しなさい。(100 字以内)\\
  \item 「確定給付企業年金に関する数理実務ガイダンス」には、
  予定利率を設定する際に検 討すべき事項として、
  事業主等へ助言を行うことが望ましいとされている点が記載されている。これを 3 つ簡潔に入力しなさい。(200 字以内)
\end{enumerate}

\end{shortanswer}

\begin{answerbox}
  \begin{enumerate}[(1)]
  \item 運用収益の長期の予測\\
  \item 直近 5 年間に発行された 10 年国債の応募者利回りの平均又は直近 1 年間 に発行された 10 年国債の応募者利回りの平均のいずれか低い率を基準と して設定されたものであること。\\
  \item 次のうちから3つを記載
  \begin{itemize}
    \item 長期的期待収益率および下限予定利率との整合性
    \item 設定した予定利率の年金財政への影響
    \item 適用後の決算時における利差損益の動向
    \item 長期的期待収益率とリスクとの整合性
    \item その他、年金財政への影響が大きいと思われる事項
  \end{itemize}
\end{enumerate}

\end{answerbox}

\newpage

\subsection{掛金の額の基準2}


\begin{InyouJobunbox}
  \begin{itemize}
    \item 法第57条
    \item 則第45条、第48条〜第49条
    \item 解釈第4-4
  \end{itemize}
\end{InyouJobunbox}


\begin{truefalse}[年金2 2019 問題1(1)\textcircled{3}]
  確定給付企業年金法施行規則第48 条第1 項第2 号の方法により積立金の額を評価している
確定給付企業年金制度において、年金資産として時価を使用するものを次の項目から全てあげなさい。

\begin{enumerate}[1. ]
  \item リスク充足額の算出
  \item 法第 62 条の規定による責任準備金の額に照らし厚生労働省令で定めるところにより
算定した額との比較による財政検証
  \item 法第 63 条の規定による最低積立基準額との比較による財政検証
  \item 法第 64 条の規定による積立上限額を超える場合の掛金の控除額の算出
\end{enumerate}
\end{truefalse}

\begin{answerbox}
  3
\end{answerbox}

\begin{fillblank}[数理人試験 2024 問題4 設問1 A,B]
  \noindent \textbf{◯DB則}
  \begin{lawstructure}
    \begin{lawarticle}{第四十五条}{掛金の額の計算に関する基準}
      （略）

      \lawparagraph{2}
      前項の標準掛金額とは、給付に要する費用（第四十三条の規定に基づき計算した通常予測給付額のうち
      計算基準日後の加入者であった期間となると見込まれる期間に係るものに限る。第二号において同じ。）
      に充てるため事業主が拠出する掛金の額であって、原則として、将来にわたって平準的に、かつ、
      加入者となる者に係る第一号の額が第二号の額を下回らないように定められる掛金の額をいう。

      \lawitem{一}
      （　　A　　）に相当する額

      \lawitem{二}
      （　　B　　）に相当する額


      \lawparagraph{3}
      （略）
      
      \lawparagraph{4}
      第一項のリスク分担型企業年金掛金額とは、給付に要する費用に充てるため事業主が拠出する額であって、第四十六条の三の規定に基づき定められる掛金の額をいう。
    \end{lawarticle}
  \end{lawstructure}
\end{fillblank}

\begin{answerbox}
    \begin{enumerate}[A: ]
    \item 標準掛金額の予想額の現価
    \item 給付に要する費用の通常の予測に基づく予想額の現価
  \end{enumerate}
\end{answerbox}


\newpage

\subsection{特別掛金}


\begin{InyouJobunbox}
  \begin{itemize}
    \item 則第46条、第112条
    \item 解釈第4-6(1)
    \item 特例的扱い第3-1(1)
  \end{itemize}
\end{InyouJobunbox}

\subsubsection{過去勤務債務の額の算定}

\begin{fillblank}[数理人試験 2021 問題1 設問1 A,B]
  \noindent \textbf{◯特例的扱い}
  \begin{lawstructure}
    \begin{lawarticle}{}{}
      \noindent 第3 財政計算時の特例
      \\

1 特別掛金の算定に係る原則的扱い
\\

(1)財政計算に用いる資産額は、基準日における純資産額(資産の評価に数理的評価を用いている場合には数理上資産額)から
（　　A　　）の額を控除した額とし、規則第四十六条第一項の規定により過去勤務債務の額を算定する場合の積立金の額には、
当該資産額を用いること。なお、規則第四 十六条第一項第一号、第二号及び第四号の方法により
特別掛金を算定する場合において、基準日以降における（　　B　　）を見込むことにより算定することができること。

    \end{lawarticle}
  \end{lawstructure}
\end{fillblank}

\begin{answerbox}
  \begin{enumerate}[A: ]
    \item 別途積立金
    \item 加入者の数または加入者の給与の額の変動
  \end{enumerate}
\end{answerbox}

\newpage

\subsubsection{特別掛金の計算方法}

\begin{fillblank}[数理人試験 H24 問題1(4)]
    \noindent \textbf{◯DB則}
  \begin{lawstructure}
    \begin{lawarticle}{第四十六条}{特別掛金額}
      前条第一項の補足掛金額のうち過去勤務債務の額（第四十三条の規定に基づき計算した通常予測給付額の現価に相当する額から標準掛金額の予想額の現価に相当する額と積立金の額を合算した額を控除した額をいう。以下同じ。）に係る掛金の額（以下「特別掛金額」という。）は、次のいずれかの方法により計算されなければならない。

      \lawitem{一}
      過去勤務債務の額を三年以上二十年以内の範囲内においてあらかじめ規約で定めた期間（以下「予定償却期間」という。）で均等に償却する方法

      \lawitem{二}
      前号の方法で計算した特別掛金額（以下この号において「下限特別掛金額」という。）及び次の表の上欄に掲げる予定償却期間ごとに同表の下欄に掲げる最短期間を予定償却期間として前号の方法で計算した特別掛金額（以下この号において「上限特別掛金額」という。）を規約で定め、併せて、毎事業年度の特別掛金額を下限特別掛金額以上、上限特別掛金額以下の範囲内において規約で定める方法
      \begin{table}[htbp]
        \centering
        \begin{tabular}{|l|l|} \hline
          予定償却期間 & 最短期間 \\ \hline
          （　　A　　）年未満 & 三年 \\ \hline
          （中略） & （中略） \\ \hline
          （　　B　　）年以上 & （　　C　　）年 \\ \hline
        \end{tabular}
      \end{table}

      \lawitem{三}
      過去勤務債務の額に百分の（　　D　　）以上百分の五十以下の範囲内において規約で定めた一定の割合を乗じて償却する方法（毎事業年度の特別掛金額を規約で定めることとし、過去勤務債務の額が当該事業年度の標準掛金額以下となるときは、当該過去勤務債務の額の全部を当該特別掛金額とすることができるものとする。）

      \lawitem{四}
      予定償却期間において、次に掲げる要件を満たすように特別掛金額を定めて償却する方法

      \hangedlawitem{イ}{3}
      特別掛金額は、過去勤務債務の額の償却開始後（　　E　　）年を経過するまでの間に（　　F　　）かつ
      引上げ額が（　　G　　）に大きくならない方法で、段階的に引き上げられるものであること。
    
      \hangedlawitem{ロ}{3}
      特別掛金額の予想額の現価に相当する額が過去勤務債務の額を下回らないこと。

      \hangedlawitem{ハ}{3}
      予定償却期間中の各期間における特別掛金額について、あらかじめ規約に定めていること。


      （以下略）
    \end{lawarticle}
  \end{lawstructure}
\end{fillblank}

\begin{answerbox}
  \begin{enumerate}[A: ]
    \item 五
    \item 十五
    \item 十
    \item 十五
    \item 五
    \item 定期的
    \item 経年的
  \end{enumerate}
\end{answerbox}

\begin{shortanswer}[数理人試験 H27 問題3 設問1]
  特別掛金額の計算方法について確定給付企業年金法施行規則第四十六条第一項に
  「過去勤務債務の額を三年以上二十年以内の範囲内においてあらかじめ規約で定めた期間で均等に償却する方法」
  (以下「元利均等償却」という。)の他3つの計算方法が規定されている。当該3つの計算方法についてその内容を簡記せよ。
\end{shortanswer}

\begin{answerbox}
  \begin{enumerate}[(1)]
    \item 「下限特別掛金額」及び「上限特別掛金額」を規約で定め、併せて、
毎事業年度の特別掛金額を下限特別掛金額以上、上限特別掛金額
以下の範囲内において規約で定める方法
    \item 過去勤務債務の額に規約で定めた一定の割合(15/100以上50/100以下)を乗じて償却する方法
    \item 過去勤務債務の額の償却開始後五年を経過するまでの間に定期的かつ
引上げ額が経年的に大きくならない方法で、段階的に引き上げられる
方法
  \end{enumerate}
\end{answerbox}

\newpage

\subsubsection{前回の過去勤務債務の償却が未完了の場合}

\begin{shortanswer}[年金2 2023 問題2(1)]
  前回の財政計算において発生した過去勤務債務の額の償却が完了していない場合の確定
給付企業年金法施行規則第４６条第 1 項第1号の方法（元利均等償却）に基づく特別掛金の
算定について、以下の（1）、（2）に答えなさい。

\begin{enumerate}[(1)]
  \item 今回の財政計算において発生した過去勤務債務の額が前回の財政計算において発生した
過去勤務債務の額のうち償却されていない額を下回るときに、特別掛金額を計算する場合
の予定償却期間の取扱いを簡潔に入力しなさい。（200 字以内）\\
\item 今回の財政計算において新たに過去勤務債務の額が発生したため、確定給付企業年金法
施行規則第４６条第２項第３号の方法（前回の財政計算において発生した過去勤務債務の
額のうち償却されていない額と今回の財政計算で新たに発生した過去勤務債務の額を合算
した額について、合理的に計算した額とする方法）で特別掛金額を計算する。前回の財政計
算において定めた予定償却期間の残存期間が三年に満たないときの予定償却期間と掛金額
の上限及び下限の取扱いを簡潔に入力しなさい。（200 字以内）
\end{enumerate}
\end{shortanswer}

\begin{answerbox}
  \begin{enumerate}[(1)]
    \item 今回の財政計算において発生した過去勤務債務の額の償却が完了する日は、前回
の財政計算において発生した過去勤務債務の額の償却が完了することとしていた
日後の日としてはならず、前回の財政計算において定めた予定償却期間の残存
期間が三年に満たないときは、予定償却期間を当該残存期間としなければなら
ない。\\
    \item 特別掛金額は、前回の財政計算において計算した特別掛金額を下回らないもの
とし、前回の財政計算において定めた特別掛金額に今回の財政計算で新たに発生
した過去勤務債務の額を三年で償却するとした場合の特別掛金額を加算した額を
上回らない範囲内で定めることができる。

この場合においては、予定償却期間を三年未満とすることができる。
  \end{enumerate}
\end{answerbox}

\begin{shortanswer}[数理人試験 H27 問題3 設問2]
  「元利均等償却」により特別掛金額を計算する場合において、前回の財政計算において定めた
  予定償却期間の残存期間が(1)3年以上のとき、(2)3年未満のときのそれぞれについて、今回の財
  政計算において発生した過去勤務債務の額の予定償却期間を設定する際の取り扱いについて簡
  記せよ。
  
  なお、今回の財政計算において発生した過去勤務債務の額は、前回の財政計算において発生した
  過去勤務債務の額のうち償却されていない額を下回っているものとする。
\end{shortanswer}

\begin{answerbox}
  \begin{enumerate}[(1)]
    \item 前回の財政計算において定めた予定償却期間の残存期間を延長してはならない。 なお、3年未満とすることはできない。
    \item 前回の財政計算において定めた予定償却期間の残存期間としなければならない。
  \end{enumerate}
\end{answerbox}

\newpage


\subsubsection{予定利率を引き下げる場合の特例}

\begin{shortanswer}[数理人試験 H27 問題3 設問3]
  今回の財政計算において繰越不足金の解消に併せて予定利率を引き下げる場合において、特別
  掛金額を「元利均等償却」により計算するときに、確定給付企業年金法施行規則第四十六条第
  五項に規定されている過去勤務債務の額の予定償却期間を設定する際の取り扱いについて簡記
  せよ。 
  
  なお、前回の財政計算において発生した過去勤務債務については既に償却が完了しているもの
  とし、今回の財政計算において予定利率以外の基礎率は変更しないものとする。
\end{shortanswer}

\begin{answerbox}
  今回の財政計算において計算した数理債務の額から、予定利率を引き下げないもの
として計算した数理債務の額を控除して得た額(以下「予定利率引下げによる過去
勤務債務の額」という。)については、予定償却期間を三年以上三十年以内の範囲
内においてあらかじめ規約で定めた期間で設定する事ができる。
今回の財政計算において計算した過去勤務債務の額から「予定利率引下げによる過
去勤務債務の額」を控除した額については予定償却期間を三年以上二十年以内の範
囲内においてあらかじめ規約で定めた期間で設定しなければならない。
\end{answerbox}


\newpage

\subsection{リスク対応掛金額}


\begin{InyouJobunbox}
  \begin{itemize}
    \item 則第46条の2
    \item 解釈第4-6(1)(2)
  \end{itemize}
\end{InyouJobunbox}

\subsubsection{リスク対応掛金の拠出方法}

\begin{fillblank}[年金2 H29 問題1(3)]
  \noindent \textbf{◯DB則}
  \begin{lawstructure}
    \begin{lawarticle}{第四十六条の二}{リスク対応掛金額}
      第四十五条第一項の補足掛金額のうち財政悪化リスク相当額に係る掛金の額（以下「リスク対応掛金額」という。）は次の各号のいずれかの方法により計算されなければならない。
    
      \lawitem{一}
      財政悪化リスク相当額から（　　A　　）（積立金の額並びに標準掛金額及び特別掛金額の予想額の現価に相当する額を
      合算した額から通常予測給付額の現価に相当する額を控除した額（当該額が零未満となる場合にあっては零とする。）をいう。）
      を控除した額（当該額が零未満となる場合にあっては零とする。）の範囲内において、あらかじめ計画的に掛金を拠出すること
      が適当であるものとして規約で定める額（以下「リスク対応額」という。）を（　　B　　）の範囲内においてあらかじめ
      規約で定めた期間（以下「予定拠出期間」という。）で均等に拠出する方法

      \lawitem{二}
      （略）

      \lawitem{三}
      リスク対応額（既にリスク対応掛金額として拠出した部分の額を除く。以下この号において同じ。）
      に百分の十五以上百分の五十以下の範囲内において規約で定めた一定の割合を乗じて拠出する方法
      （毎事業年度のリスク対応掛金額を規約で定めることとし、リスク対応額が当該事業年度の（　　C　　）
      となるときは、当該リスク対応額の全部をリスク対応掛金額とすることができるものとする。）

      \lawitem{四}
      予定拠出期間において、次に掲げる要件を満たすようにリスク対応掛金額を定めて拠出する方法

      \hangedlawitem{イ}{3}
      リスク対応掛金額は、拠出開始後五年を経過するまでの間に定期的かつ引上げ額が（　　D　　）方法で、段階的に引き上げられるものであること。

      \hangedlawitem{ロ}{3}
      リスク対応掛金額の予想額の現価に相当する額がリスク対応額を上回らないこと。

      \hangedlawitem{ハ}{3}
      予定拠出期間中の各期間におけるリスク対応掛金額について、あらかじめ規約に定めていること。

      （略）


      \lawparagraph{4}
      特別掛金額の予定償却期間の残存期間はリスク対応掛金額の予定拠出期間の残存期間より（　　E　　）でなければならない。
    \end{lawarticle}
  \end{lawstructure}
\end{fillblank}

\begin{answerbox}
  \begin{enumerate}[A: ]
    \item 対応前リスク充足額
    \item 五年以上二十年以下
    \item 標準掛金額以下
    \item 経年的に大きくならない方法
    \item 短い期間
  \end{enumerate}
\end{answerbox}

\begin{fillblank}[数理人試験 2020 問題1 設問1 A,B]
  \noindent \textbf{◯DB則}
  \begin{lawstructure}
    \begin{lawarticle}{第四十六条の二}{リスク対応掛金額}
      第四十五条第一項の補足掛金額のうち財政悪化リスク相当額に係る掛金の額（以下「リスク対応掛金額」という。）は次の各号のいずれかの方法により計算されなければならない。

      \lawitem{一}
      （略）

      \lawitem{二}
      前号の方法で計算したリスク対応掛金額（以下この号において「下限リスク対応掛金額」という。）及び次の表の上欄に掲げる予定拠出期間ごとに同表の下欄に掲げる最短期間を予定拠出期間として前号の方法で計算したリスク対応掛金額（以下この号において「上限リスク対応掛金額」という。）を規約で定め、併せて、毎事業年度のリスク対応掛金額を下限リスク対応掛金額以上、上限リスク対応掛金額以下の範囲内において規約で定める方法
      \begin{table}[htbp]
        \centering
        \begin{tabular}{|l|l|} \hline
          予定償却期間 & 最短期間 \\ \hline
          （　　A　　） & 五年 \\ \hline
          （略） & （略） \\ \hline
          （　　B　　） & 七年 \\ \hline
          （略） & （略） \\ \hline
        \end{tabular}
      \end{table}

      （以下略）
    \end{lawarticle}
  \end{lawstructure}
\end{fillblank}


\begin{answerbox}
  \begin{enumerate}[A: ]
    \item 九年未満
    \item 十一年以上十三年未満
  \end{enumerate}
\end{answerbox}

\newpage

\subsubsection{リスク対応掛金の変更・終了}

\begin{fillblank}[数理人試験 2020 問題1 設問1 C〜G]
  \noindent \textbf{◯DB則}
  \begin{lawstructure}
    \begin{lawarticle}{第四十六条の二}{リスク対応掛金額}
      （略）

      \lawparagraph{2}
      リスク対応掛金額の拠出が完了していない場合であって、次の各号に掲げる場合に該当することとなったときには、当該各号に定めるところによりリスク対応掛金額を変更することができる。

      \lawitem{一}
      財政計算を行い、新たに過去勤務債務の額が発生する場合　増加する特別掛金額の予想額の現価に相当する額が
      リスク対応掛金額の予想額の現価に相当する額の（　　A　　）範囲内でリスク対応掛金額を減少させること。

      \lawitem{二}
      第五十条各号に掲げる場合（同条（　　B　　）に掲げる場合を除く。）　前項の規定に従い、リスク対応掛金額を計算すること。

      \lawitem{三}
      法第五十八条第一項の規定に基づく財政再計算において、財政悪化リスク相当額から（　　C　　）
      （積立金の額と標準掛金額、特別掛金額及び当該財政再計算による変更前のリスク対応掛金額の予想額の現価を
      合算した額から通常予測給付額の現価に相当する額を控除した額（当該額が零未満となる場合にあっては零とする。）
      をいう。次項において同じ。）を控除した額（当該額が零未満となる場合にあっては零とする。）が、
      前項の規定に基づきリスク対応掛金額を計算したとき（リスク対応掛金額を変更した場合にあっては、
      当該変更のうちの直前の変更をしたとき）から（　　D　　）する場合　当該（　　D　　）した額を上回らない範囲で同項第一号
      のリスク対応額を定め、同項の規定に基づき計算したリスク対応掛金額に相当する額を変更前のリスク対応掛金額
      に加算すること。

      \lawparagraph{3}
      法第五十八条第一項の規定に基づく財政再計算において、
      （　　C　　）が財政悪化リスク相当額を上回ることとなる場合には、
      上回らないようにリスク対応掛金額を減少させ、又はリスク対応掛金額の（　　E　　）しなければならない。

      \lawparagraph{4}
      （略）
    \end{lawarticle}
  \end{lawstructure}
\end{fillblank}

\begin{answerbox}
  \begin{enumerate}[A: ]
    \item 減少額を下回らない
    \item 第四号のニ
    \item 対応後リスク充足額
    \item 増加
    \item 拠出を終了
  \end{enumerate}
\end{answerbox}

\begin{shortanswer}[年金2 2023 問題2(3)]
  リスク対応掛金を拠出しているＡ社から、拠出が完了する前であるがリスク対応掛金額の
変更を検討したいと申し出があった。加えて、「その他当該確定給付企業年金に係る事情に
著しい変動があった場合（確定給付企業年金法施行規則第５０条第４号ホ）」にリスク対応
掛金額の変更が可能とされているが、これはどのような場合を指すのか問い合わせを受け
た。これに該当する例示を５つ簡潔に入力しなさい。（250字以内）

なお、記載にあたっては、「確定給付企業年金に関する数理実務ガイダンス」に掲げる事
例の中から記載すること。
\end{shortanswer}

\begin{answerbox}
  以下のうち５つを記載。
  \begin{enumerate}
    \item 財政悪化リスク相当額の算定方法を変更する場合。
    \item 繰越不足金の解消を目的として実施する財政再計算。
    \item 給付支払いのための積立金が不足することにより、規則第64条の掛金（特別掛金）を拠出した場合に実施する財政再計算。
    \item 特別掛金の引き下げもしくは廃止を目的として別途積立金を取り崩して実施する財政再計算。
    \item 資産の評価方法の変更により実施する財政再計算
    \item 予定利率を変更することを目的として実施する財政再計算
    \item リスク対応掛金を新たに設定することを目的として実施する財政再計算
    \item リスク算定告示別表に定められているリスク係数対象資産別のリスク係数が変更された場合
    \item 資産額の大幅な変動
    \item 政策的資産構成割合または実際の資産構成割合が大幅に変更された場合
    \item 政策的資産構成割合と実際の資産構成割合の乖離が大きい場合
  \end{enumerate}
\end{answerbox}

\begin{shortanswer}
  確定給付企業年金関連法令および関連通知の内容を踏まえて、次の(1)、(2)について簡記しな
さい。なお、リスク分担型企業年金でない確定給付企業年金制度を想定して記載すること。\\

\begin{enumerate}[(1)]
  \item リスク対応掛金額を変更できる場合の要件と変更内容\\
  \item リスク対応掛金額を変更しなければならない場合の要件と変更内容
\end{enumerate}
\end{shortanswer}

\begin{answerbox}
  \begin{enumerate}[(1)]
    \item \begin{itemize}
      \item 財政計算を行い、新たに過去勤務債務が発生する場合\\
増加する特別掛金収入現価の範囲内でリスク対応掛金額を減少させることができる。
\\
      \item 少なくとも五年ごとに実施する財政再計算において、財政悪化リスク相当額から対応後リスク充足額を控除した額が、前回財政計算時より増加する場合\\
増加した額を上回らない範囲でリスク対応額を定め、それに応じたリスク対応掛金額を変更前のリスク対応額に加算することができる。
\\
      \item 少なくとも五年ごとに実施する財政再計算以外の財政再計算に該当する場合（過去勤務債務の額の予定償却期間を短縮しようとする場合または定率償却の償却割合を増加させようとする場合を除く）\\
リスク対応掛金額をあらためて計算することができる。\\
    \end{itemize}
    
    \item 少なくとも五年ごとに実施する財政再計算において、対応後リスク充足額が財政悪化リスク相当額を上回ることとなる場合\\
上回らないようにリスク対応額を減少させ、またはリスク対応掛金額の拠出を終了しなければならない。
  \end{enumerate}
\end{answerbox}

\begin{shortanswer}[年金2 2025 問題2(1)]
  リスク対応掛金額の変更について以下の(A)、(B)に答えなさい。
\\
  \begin{enumerate}[(A)]
    \item 確定給付企業年金法施行規則第 46 条の 2 第 2 項に定める、リスク対応掛金額の拠出が完了する前にリスク対応掛金額を変更することができるとされる次の(1)及び(2)の場合 について、リスク対応掛金額の変更方法をそれぞれ簡潔に入力しなさい。(各 150 字以 内)
    \\
    \begin{enumerate}[(1)]
      \item 財政計算を行い、新たに過去勤務債務の額が発生する場合\\
      \item 確定給付企業年金法第 58 条第 1 項の規定に基づく財政再計算において、財政悪化 リスク相当額から対応後リスク充足額(※1)を控除した額(当該額が零未満とな る場合にあっては零とする。)が、リスク対応掛金額を計算したとき(※2)から 増加する場合\\
      \\
      (※1)対応後リスク充足額とは、積立金の額と標準掛金額、特別掛金額及び当該財政 再計算による変更前のリスク対応掛金額の予想額の現価を合算した額から通常 予測給付額の現価に相当する額を控除した額(当該額が零未満となる場合にあ っては零とする。)をいう。\\
      (※2)直前の財政計算において初めてリスク対応掛金を算定したものとする。\\
    \end{enumerate}
    \item 確定給付企業年金法施行規則第 46 条の 2 第 3 項に定める、リスク対応掛金額を減少さ せ、又はリスク対応掛金額の拠出を終了しなければならない場合を簡潔に入力しなさ い。(100 字以内)
  \end{enumerate}
\end{shortanswer}

\begin{answerbox}
  \begin{enumerate}[(A)]
    \item \begin{enumerate}[(1)]
      \item 増加する特別掛金額の予想額の現価に相当する額がリスク対応掛金
額の予想額の現価に相当する額の減少額を下回らない範囲内でリス
ク対応掛金額を減少させる\\
      \item 「財政悪化リスク相当額から対応後リスク充足額を控除した額が増
加した額」を上回らない範囲でリスク対応額を定め、算定したリス
ク対応掛金額に相当する額を変更前のリスク対応掛金額に加算する\\
    \end{enumerate}
    \item 法第 58 条第 1 項の規定に基づく財政再計算において、対応後リスク充足 額が財政悪化リスク相当額を上回ることとなる場合
  \end{enumerate}
\end{answerbox}

\newpage

\subsection{リスク分担型企業年金掛金額}

\begin{InyouJobunbox}
  \begin{itemize}
    \item 則第46条の3
    \item 解釈第4-6(3)
    \item 特例的扱い第6
  \end{itemize}
\end{InyouJobunbox}


\begin{truefalse}[年金2 2021 問題1(1)\textcircled{1}]
  「確定給付企業年金に関する数理実務ガイダンス」に記載されているリスク分担型企業年
金の取り扱いとして適当なものをすべて選びなさい。
\begin{enumerate}
  \item 財政方式によらず、通常予測給付現価の算定対象に将来加入者を含むことを基本とす
る。
  \item 非継続基準は結果的に1.00 となるが、最低積立基準額は計算する必要がある。
  \item 積立上限額の計算を要しない。
  \item 承継事業所償却積立金は設定できない。
\end{enumerate}

\end{truefalse}

\begin{answerbox}
  1,2,3,4
\end{answerbox}

\begin{explanationbox}
1: DBガイダンス第4節1(82p)

2: DBガイダンス第4節2(82p)

3: DBガイダンス第4節3(83p)

4: DBガイダンス第4節6(83p)
\end{explanationbox}

\begin{shortanswer}[年金2 2020 問題2(2)]
  リスク分担型企業年金において、財政悪化リスク相当額の価格変動リスクにおける「定常状態の積立金の額の予想額」の算定方法を 3 つ簡記しなさい。
\end{shortanswer}

\begin{answerbox}
  \begin{itemize}
    \item 「定常状態の通常予測給付額の現価−定常状態のリスク分担型年金掛金額現価」とする方法
    \item 「定常状態の通常予測給付額の現価−定常状態のリスク分担型年金掛金額現価+計算基準日における超過財源」とする方法
    \item 「計算基準日における超過財源を0とするような給付調整を恒久的に行うことを仮定した定常状態の通常予測給付額の現価−定常状態のリスク分担型企業年金掛金額現価」とする方法
  \end{itemize}
\end{answerbox}

\begin{explanationbox}
  DBガイダンス補足事項（財政悪化リスク相当額）【例示１】（p.131）

DBガイダンス補足事項（財政悪化リスク相当額）【例示２】（p.134）

DBガイダンス補足事項（財政悪化リスク相当額）【例示３】（p.135）
\end{explanationbox}

\begin{fillblank}[数理人試験 2024 問題4 設問2 C,D]
  \noindent \textbf{◯DB則}
  \begin{lawstructure}
    \begin{lawarticle}{第四十六条の三}{リスク分担型企業年金掛金額}
      リスク分担型企業年金を実施するとき又はリスク分担型企業年金を実施している場合であって
      （　　A　　）（掛金の額に係る規約の変更を行う場合に限る。）におけるリスク分担型企業年金掛金額は、
      当該リスク分担型企業年金の掛金の額を第四十五条第一項の標準掛金額、補足掛金額その他の掛金の額に
      区分して定めることとしたならば当該実施又は当該変更による財政計算において計算されることとなる標準掛金額と
      補足掛金額とを合算した額とする方法により計算されなければならない。

      \lawparagraph{2}
      （略）

      \lawitem{一}
      （略）

      \lawitem{二}
      当該再計算において（　　B　　）として規約で定める額を前条第一項第一号の
      リスク対応額とみなして同号の方法により計算した額を追加して拠出する方法

      \lawitem{三}
      （略）
    \end{lawarticle}
  \end{lawstructure}  
\end{fillblank}

\begin{answerbox}
  \begin{enumerate}[A: ]
    \item 給付の設定を変更するとき
    \item 計画的に掛金を拠出することが適当である額
  \end{enumerate}
\end{answerbox}

\newpage

\subsection{次回財政計算までに発生する積立不足の予想額とその償却}

\begin{InyouJobunbox}
  \begin{itemize}
    \item 則第44条、第47条
  \end{itemize}
\end{InyouJobunbox}

\begin{fillblank}[数理人試験 2025 問題1 設問2 E〜G]
  \noindent \textbf{◯DB則}
  \begin{lawstructure}
    \begin{lawarticle}{第四十四条}{次回の財政再計算までに発生する積立不足の予想額}
      前条の規定に基づき掛金の額を計算する場合において、次に掲げる事情によって、
      次回の財政再計算までの間に積立金の額が法第六十条第二項に規定する責任準備金の額
      （以下「責任準備金の額」という。）又は同条第三項に規定する最低積立基準額（以下「最低積立基準額」という。）
      を下回ることが予想される場合にあっては、当該下回ることが予想される額のうちいずれか（　　A　　）を
      前条の規定に基づき計算した（　　B　　）に相当する額に加算することができる。

      \lawitem{一}
      （　　C　　）が前条第二項第一号の予定利率よりも低いこと。

      \lawitem{二}
      加入者の数が一時的に著しく変動することが見込まれること。

      \lawitem{三}
      加入者の給与の額その他これに類するものが一時的に著しく変動することが見込まれること。
    \end{lawarticle}
  \end{lawstructure}
\end{fillblank}

\begin{answerbox}
  \begin{enumerate}[A: ]
    \item 大きい額の現価
    \item 通常予測給付額の現価
    \item 積立金の運用利回りの予測
  \end{enumerate}
\end{answerbox}

\begin{shortanswer}[年金1 H27 問題2(2)]
  次回の財政再計算までに発生する積立不足の予想額について、確定給付企業年金法施行規則第
44 条に記載されている内容を簡記しなさい。なお、同条第 1 号から第 3 号に掲げられる３つの事
情についても記載すること。
\end{shortanswer}

\begin{answerbox}
  掛金の額を計算する場合において、次に掲げる事情によって、
  次回の財政再計算までの間に積立金の額が責任準備金の額または最低積立基準額を下回ることが予想される場合にあっては、
  当該下回ることが予想される額のうちいずれか大きい額の現価を給付に要する費用の予想額の現価に加算することができる。
  \begin{enumerate}
    \item 積立金の運用利回りの予測が予定利率よりも低いこと
    \item 加入者の数が一時的に著しく変動することが見込まれること
    \item 加入者の給与の額その他これに類するものが一時的に著しく変動することが見込まれること
  \end{enumerate}
\end{answerbox}

\newpage

\subsection{簡易な基準に基づくDB}

\begin{InyouJobunbox}
  \begin{itemize}
    \item 則第52条
  \end{itemize}
\end{InyouJobunbox}

\begin{fillblank}[年金1 H27 問題1(3)]
  \noindent \textbf{◯DB則}
  \begin{lawstructure}
    \begin{lawarticle}{第五十二条}{簡易な基準に基づく確定給付企業年金の掛金の額の算定}
      計算基準日における加入者の数が（　　A　　）に満たない確定給付企業年金（（　　B　　）を除く。）
      の掛金の額は、
      第四十三条の規定にかかわらず、次に定めるところにより計算することができる。

      \lawitem{一}
      基礎率のうち予定利率及び予定死亡率のみを用いること。ただし、給付の額が令第二十四条第一項第三号の方法により計算される場合（第二十五条の規定により令第二十四条第一項第三号の方法を組み合わせている場合を含む。）にあっては、同号の再評価に用いる指標の予測を用いること。

      \lawitem{二}
      予定利率は、（　　C　　）以上四・〇パーセント以下の範囲内とすること。

      \lawitem{三}
      予定死亡率は、第六十二条第一号ロに規定する予定死亡率とすること。

      \lawitem{四}
      令第二十四条第三項の（　　D　　）を行わないこと。

      \lawitem{五}
      障害給付金を支給しないこと。

      \lawitem{六}
      遺族給付金を支給する場合にあっては、当該遺族給付金の額は、老齢給付金の保証期間の残存期間について支給する給付の現価に相当する金額又は脱退一時金（法第二十九条第一項第二号に規定する脱退一時金をいう。以下同じ。）の額以下となっていること。

      \lawparagraph{2}
      （　　B　　）（（　　E　　）を除く。）の掛金の額は、第四十三条の規定にかかわらず、契約者価額の計算に用いる予定利率及び予定死亡率を用い、前項第一号、第五号及び第六号に規定するところにより計算することができる。
    
      （以下略）
    \end{lawarticle}
  \end{lawstructure}
\end{fillblank}

\begin{answerbox}
  \begin{enumerate}[A: ]
    \item 五百人
    \item 受託保証型確定給付企業年金
    \item 下限予定利率
    \item 給付の額の改定
    \item 閉鎖型受託保証型確定給付企業年金
  \end{enumerate}
\end{answerbox}

\begin{truefalse}[年金2 2020 問題1(1)\textcircled{5}]
  確定給付企業年金法施行規則第 52 条第 4 項では、事業主等が規約の変更を行い、
  受託保証型確定給付企業年金を実施する場合には、数理債務の額から契約者価額を控除した額を\underline{特例掛金額
  として一括して拠出しなければならない}とされている。
\end{truefalse}

\begin{answerbox}
  \textbf{誤}

  特別掛金額として一括して拠出することができる
\end{answerbox}

\begin{truefalse}[年金2 2024 問題1(1)ウ]
  簡易な基準に基づく確定給付企業年金(受託保証型確定給付企業年金を除く)として取り扱うために必要な事項について、適当なものをすべて選びなさい。

\begin{enumerate}
  \item 計算基準日における加入者の数が 300 人未満であること。
  \item 予定利率は、下限予定利率以上 3.0 パーセント以下の範囲内とすること。
  \item 掛金の額は定額法(加入者数に比例して定める掛金)により算定すること。
  \item 財政悪化リスク相当額を零とすること。
\end{enumerate}
\end{truefalse}

\begin{answerbox}
  4
\end{answerbox}

\begin{explanationbox}
  \begin{enumerate}
    \item ×300人未満→◯500人未満（則第52条第1項）
    \item ×3.0パーセント以下→◯4.0パーセント以下（則第52条第1項第2号）
    \item そのような規定はない。ちなみに、簡易DC（※廃止済）の掛金の決め方は定額のみ。
    \item 正しい。（リスク算定告示第7条）
  \end{enumerate}
\end{explanationbox}

\begin{fillblank}[数理人試験 H25 問題1(3)e~i 改題]
  \noindent \textbf{◯DB則}
  \begin{lawstructure}
    \begin{lawarticle}{第五十二条}{簡易な基準に基づく確定給付企業年金の掛金の額の算定}
      計算基準日における加入者の数が（　　A　　）に満たない確定給付企業年金（（　　B　　）を除く。）
      の掛金の額は、
      第四十三条の規定にかかわらず、次に定めるところにより計算することができる。

      \lawitem{一}
      基礎率のうち予定利率及び予定死亡率のみを用いること。ただし、給付の額が令第二十四条第一項第三号の方法により計算される場合（第二十五条の規定により令第二十四条第一項第三号の方法を組み合わせている場合を含む。）にあっては、同号の再評価に用いる指標の予測を用いること。

      \lawitem{二}
      予定利率は、（　　C　　）以上（　　D　　）以下の範囲内とすること。

      \lawitem{三}
      予定死亡率は、第六十二条第一号ロに規定する予定死亡率とすること。

      \lawitem{四}
      令第二十四条第三項の給付の額の改定を行わないこと。

      \lawitem{五}
      （　　E　　）を支給しないこと。

      \lawitem{六}
      遺族給付金を支給する場合にあっては、当該遺族給付金の額は、老齢給付金の保証期間の残存期間について支給する給付の現価に相当する金額又は脱退一時金（法第二十九条第一項第二号に規定する脱退一時金をいう。以下同じ。）の額以下となっていること。

      （以下略）
    \end{lawarticle}
  \end{lawstructure}
\end{fillblank}

\begin{answerbox}
  \begin{enumerate}[A: ]
    \item 五百人
    \item 受託保証型確定給付企業年金
    \item 下限予定利率
    \item 四・〇パーセント
    \item 障害給付金
  \end{enumerate}
\end{answerbox}

\begin{truefalse}[年金2 2025 問題1(1)ウ\ajRoman{4}]
  以下の文章の正誤を判定し、誤っている場合は正しい内容に改めなさい。
\\

  簡易な基準に基づく確定給付企業年金の財政計算において基準死亡率に乗じる率 は、加入者については 0 である。
\end{truefalse}

\begin{answerbox}
  \textbf{正}
\end{answerbox}

\newpage

\subsection{財政再計算}

\begin{InyouJobunbox}
  \begin{itemize}
    \item 法第58条
    \item 則第50条〜第51条
  \end{itemize}
\end{InyouJobunbox}

\begin{fillblank}[数理人試験 H26 問題1(3)h]
  \noindent \textbf{◯DB法}
  \begin{lawstructure}
    \begin{lawarticle}{第五十八条}{財政再計算}
      事業主等は、少なくとも五年ごとに前条の基準に従って掛金の額を再計算しなければならない。

      \lawparagraph{2}
      事業主等は、前項の規定にかかわらず、（　　A　　）が著しく変動した場合その他の厚生労働省令で定める場合は、前条の基準に従って、速やかに、掛金の額を再計算しなければならない。
    \end{lawarticle}
  \end{lawstructure}
\end{fillblank}

\begin{answerbox}
  \begin{enumerate}[A: ]
    \item 加入者の数
  \end{enumerate}
\end{answerbox}

\begin{fillblank}[数理人試験 H27 問題1 設問4]
  \noindent \textbf{◯DB則}
  \begin{lawstructure}
    \begin{lawarticle}{第五十条}{財政再計算を行う場合}
      法第五十八条第二項の厚生労働省令で定める場合は、次のとおりとする。

      \lawitem{一}
      法第七十六条第一項の規定により基金を合併する場合（同条第三項の規定により合併により基金を設立する場合を除く。）

      \lawitem{二}
      法第七十七条第一項の規定により基金を分割する場合（同条第四項の規定により分割により基金を設立する場合を除く。）

      \lawitem{三}
      法第八十条第二項又は法第八十一条第二項の規定により加入者等に係る給付の支給に関する権利義務を承継する場合（新たに規約型企業年金を実施することとなる場合又は新たに基金を設立することとなる場合を除く。）

      \lawitem{四}
      次に掲げる場合（掛金の額に係る規約の変更を行う必要がない場合を除く。）

      \hangedlawitem{イ}{3}
      （　　A　　）が前回の（　　B　　）の計算基準日における（　　A　　）に比べて著しく増加又は減少した場合

      \hangedlawitem{ロ}{3}
      加入者の資格又は（　　C　　）を変更する場合

      \hangedlawitem{ハ}{3}
      法第七十九条第一項又は第二項の規定により加入者等に係る給付の支給に関する権利義務を（　　D　　）する場合

      \hangedlawitem{ニ}{3}
      過去勤務債務の額の予定償却期間を（　　E　　）しようとする場合又は第四十六条第一項第三号の一定の割合を（　　F　　）させようとする場合

      \hangedlawitem{ホ}{3}
      その他当該確定給付企業年金に係る事情に著しい変動があった場合
    \end{lawarticle}
  \end{lawstructure}
\end{fillblank}

\begin{answerbox}
  \begin{enumerate}[A: ]
    \item 加入者の数
    \item 財政計算
    \item 給付の設計
    \item 移転または承継
    \item 短縮
    \item 増加
  \end{enumerate}
\end{answerbox}

\begin{shortanswer}[年金2 2021 問題2(1)]
  確定給付企業年金法施行規則第50 条第4 号イにおいて加入者数の大幅変動は財政再計算の
該当事由になると定められているが、「確定給付企業年金に関する数理実務ガイダンス」に記
載されている、掛金率の洗い替えや基礎率の洗い替えを行わなくても良いとする例示を 3 つ
簡記しなさい。
\end{shortanswer}

\begin{answerbox}
  \begin{itemize}
    \item 基礎率を見直すことにより、掛金率が低下する場合
    \item 次回再計算まで別途積立金が留保できる場合
    \item 掛金率の上昇がわずかであり、次回財政再計算で掛金の手当が可能である場合
  \end{itemize}
\end{answerbox}

\begin{explanationbox}
  DBガイダンス第５節1[留意事項](2)。(p.87)
\end{explanationbox}

\begin{shortanswer}[数理人試験 2023 問題4 設問2]
  確定給付企業年金法第58条第２項における「厚生労働省令で定める場合」として確定給付企
業年金法施行規則第50条に規定されている財政再計算を実施する事由を「基金を合併する場
合」「基金を分割する場合」「規約型企業年金から基金、又は、基金から規約型企業年金への
移行」以外で２つ記載せよ。
\end{shortanswer}

\begin{answerbox}
  次のうち２点を記載すれば良い。
  \begin{itemize}
    \item 加入者の数が前回の財政計算の計算基準日における加入者の数に比べて著しく増加または減少した場合
    \item 加入者の資格または給付の設計を変更する場合
    \item 実施事業所の全部または一部における給付の支給に関する権利義務を承継または移転する場合
    \item 過去勤務債務の額の予定償却期間の短縮または定率償却の償却割合の増加を行う場合
    \item その他当該確定給付企業年金にかかる事情に著しい変動があった場合
  \end{itemize}
\end{answerbox}

\begin{truefalse}[年金2 2025 問題1(1)ア\ajRoman{4}]
以下の文章の正誤を判定し、誤っている場合は正しい内容に改めなさい。
\\

  加入者の数が前事業年度の末日における加入者の数に比べて著しく増加又は減少
  した場合は掛金の額を再計算しなければならない。ただし、財政の健全性が確保
  できる場合においては掛金率の洗い替え、基礎率の洗い替えを行わなくても良い。
\end{truefalse}

\begin{answerbox}
  \textbf{誤}

  「前事業年度の末日」ではなく「直前の財政計算の計算基準日」
\end{answerbox}

\begin{truefalse}[数理人試験 H15 問題1(5)]
  確定給付企業年金において、次の事由A〜Eのうち財政再計算に該当しないものをすべて選びなさい。
  なお、事由A〜Eのいずれも掛金の額に係る規約の変更が必要な場合であるものとする。

  \begin{enumerate}[A. ]
    \item 加入者の資格又は給付の設計を変更する場合
    \item 決算において積立金の額が責任準備金の額から許容繰越不足金を控除した額を下回った場合
    \item 決算において積立金の額が最低積立基準額の90\%を下回った場合
    \item 決算において積立金の額が積立上限額を上回った場合
    \item 企業年金基金を合併する場合(合併により企業年金基金を設立する場合は除く。)
  \end{enumerate}

\end{truefalse}

\begin{answerbox}
  C,D
\end{answerbox}

\newpage

\subsection{リスク算定告示：財政悪化リスク相当額の算定}

\begin{InyouJobunbox}
  \begin{itemize}
    \item リスク算定告示第1条〜第2条、第5条〜第7条、別表
    \item 承認認可基準（別紙１の２）(7)
  \end{itemize}
\end{InyouJobunbox}

\begin{truefalse}[年金2 2018 問題1(1)\textcircled{3}]
以下の文章の正誤を判定し、誤っている場合は正しい内容に改めなさい。
\\

  確定給付企業年金において、平成30 年1 月 1 日以降を計算基準日として行う\underline{財政決算}毎に
適切な算定方法によって財政悪化リスク相当額を算定する。（ただし、簡易な基準を使用し
ている場合は財政悪化リスク相当額を零とする。）
\end{truefalse}

\begin{answerbox}
  \textbf{誤}

  財政計算
\end{answerbox}

\begin{explanationbox}
  DBガイダンス第5節1（p.86）
\end{explanationbox}

\begin{shortanswer}[年金2 2018 問題2(1)\textcircled{1}]
  次の(1)、(2)について特別算定方法により財政悪化リスク相当額を算定しなければなら
ない場合について簡記しなさい。

\begin{enumerate}[(1)]
  \item リスク分担型企業年金でない確定給付企業年金
  \item リスク分担型企業年金
\end{enumerate}
\end{shortanswer}

\begin{answerbox}
  \begin{enumerate}[(1)]
    \item 積立金の資産に占める「国内債券・国内株式・外国債券・外国株式・一般勘定・短期資産」以外の資産の割合が20パーセント以上である場合。
    \item \begin{itemize}
      \item リスク算定用資産構成割合において、「国内債券・国内株式・外国債券・外国株式・一般勘定・短期資産」以外の資産構成割合が10パーセント以上である場合。
      \item 年金数理人から、年金数理に関する業務に係る書類に予定利率以外の基礎率の変動を勘案すべき旨の所見が付された場合。
    \end{itemize}
  \end{enumerate}
\end{answerbox}

\begin{explanationbox}
  リスク算定告示第2条第2項。
\end{explanationbox}


\newpage

\subsection{リスク算定告示：特別算定方法}


\begin{InyouJobunbox}
  \begin{itemize}
    \item リスク算定告示第3条〜第4条
    \item 承認認可基準（別紙１の２）(1)〜(6)
  \end{itemize}
\end{InyouJobunbox}

\subsubsection{承認不要な特別算定方法}

\begin{truefalse}[年金2 2022 問題1(1)A]
  リスク分担型企業年金以外の確定給付企業年金において、 「確定給付企業年金法施行規
則第四十三条第一項に規定する通常の予測を超えて財政の安定が損なわれる危険に対応す
る額の算定方法」第３条第１項各号に定める方法（厚生労働大臣の承認を必要としない特
別算定方法）の取扱いとして適当なものをすべて挙げなさい。

\begin{enumerate}[1. ]
  \item 価格変動リスクを計算基準日時点の資産ではなく政策的資産構成割合に基づき標準算定
方法に準じて算定する方法について、政策的資産構成割合のその他の資産が1 割以上の
場合は特別算定方法の承認が必要である。
  \item 価格変動リスクを計算基準日時点の資産ではなく政策的資産構成割合に基づき標準算定
方法に準じて算定する方法について、政策的資産構成割合の見直しを予定している場合
にあって、見直し後の政策的資産構成割合を用いる場合は特別算定方法の承認が必要で
ある。
  \item 運用環境の変化により資産の価格が変動する以外の理由により積立金の額に増減が生じ
る場合に、価格変動リスクをその積立金の増減を勘案して財政悪化リスク相当額を算定
する方法としては、積立金の増減を見込まずに標準算定方法や1の方法で算定した財政
悪化リスク相当額に、増減前後の積立金の比率を乗じる方法が含まれる。
  \item 負債変動リスクを予定利率が１％（下限予定利率を下回る場合は、下限予定利率までの
低下幅）低下した場合の数理債務の増加額として算定し、価格変動リスク（標準算定方
法又は1、3の方法で算定したもの）に当該増加額を加算する方法は特別算定方法の承
認は不要であるが、予定利率の低下幅1％未満（下限予定利率を下回ることを理由に１％
未満となる場合を除く）を使用する場合は、合理的な理由を付すことで特別算定方法の
承認を不要とできる。
\end{enumerate}
\end{truefalse}

\begin{answerbox}
  3
\end{answerbox}

\begin{explanationbox}
  1. ×1割以上→◯2割以上（リスク算定告示第3条第1項第1号イ、承認認可基準(別紙1の2)(3)）

  2. 政策的資産構成割合の見直しを予定している場合は、見直し後の資産構成割合を用いることも可能。（承認認可基準(別紙1の2)(3)）

  3. 記載の通り。（DBガイダンス補足事項(リスク分担型以外)4(2) p.111）

  4. 予定利率の低下幅を1\%（下限予定利率を下回る場合は、下限予定利率までの低下幅）以外とする場合には、特別算定方法の承認が必要。（DBガイダンス補足事項(リスク分担型以外)4(3) p.112）
\end{explanationbox}

\begin{shortanswer}
  特別算定方法のうち、第三条第一項に規定される厚生労働大臣の承認が不要となる場合につい
て簡記せよ。
\end{shortanswer}

\begin{answerbox}
  リスク分担型企業年金でない確定給付企業年金であって、次に掲げる算定方法で財政悪化リスク相当額を算定する場合（ただし、そのほかの資産の割合が20\%以上である場合を除く）。
  \begin{enumerate}[(1)]
    \item 計算基準日における積立金の額にリスク算定用資産構成割合を乗じて得た額を用いて標準算定方法に準じて算定する方法
    \item 計算基準日以後に積立金、脱退一時金額、解約手当金相当額、残余財産もしくは個人別管理資産の移換を受ける場合、積立金もしくは脱退一時金相当額を移換する場合または実施事業所が増加もしくは脱退一時金相当額を移換する場合または実施事業所が増加もしくは減少する場合に、計算基準日における積立金の額に当該増減を反映して標準算定方法に準じて算定する方法
    \item 財政悪化リスク相当額を次の合計とする場合
    \begin{enumerate}[(A)]
      \item 価格変動リスクとして、標準算定方法または上記(1)、(2)の方法で算定した額
      \item 負債変動リスクとして、予定利率を１％引き下げた場合の「通常予測給付現価ー掛金収入現価」の増加額（１％引き下げた場合の予定利率が下限予定利率を下回る場合は、当該率を予定利率とみなして算定したことによる増加額）
    \end{enumerate}
  \end{enumerate}
\end{answerbox}

\newpage

\subsubsection{特別算定方法の承認}

\begin{shortanswer}[年金2 2018 問題2(1)\textcircled{2}]
  厚生労働大臣が特別算定承認をするための要件を3点簡記しなさい。
\end{shortanswer}

\begin{answerbox}
  \begin{itemize}
    \item 財政悪化リスク相当額を、通常予測給付額の現価に相当する額（リスク分担型企業年金の場合は、調整前給付現価相当額）から掛金の額（標準掛金額と特別掛金額を合算した額またはリスク分担型企業年金掛金額）の予想額の現価に相当する額と積立金の額を合算した額を控除した額の二十年に一回の頻度で発生すると予想される最大額とするものであること。
    \item 資産の価格変動により積立金の額が低下する危険を考慮するものであり、かつ、基礎率と実績とが乖離することに伴い発生しうる危険を考慮するよう努めているものであること。ただし、リスク分担型企業年金を実施する場合にあっては、予定利率と実績とが乖離することに伴い発生しうる危険を考慮するものであること。
    \item 信頼できるデータ、情報及び手法を用いるものであること。
  \end{itemize}
\end{answerbox}

\begin{explanationbox}
  リスク算定告示第3条第4項。
\end{explanationbox}

\begin{fillblank}[]
  \noindent \textbf{◯算定告示}
  \begin{lawstructure}
    \begin{lawarticle}{第三条}{特別算定方法の承認}
      （略）

      \lawparagraph{4}
      厚生労働大臣は、特別算定承認の申請があった場合において、特別算定方法が次に掲げる要件に該当すると認めるときは、
      特別算定承認をするものとする。

      \lawitem{一}
      財政悪化リスク相当額を、通常予測給付額の現価に相当する額(リスク分担型企業年金の場合にあっては、
      調整前給付現価相当額)から掛金の額(規則第四十五条第二項に規定する標準掛金額(規則第四十六条第一項
      に規定する過去勤務債務の額がある場合にあっては、同項に規定する特別掛金額を含むことができる。)
      又はリスク分担型企業年金掛金額をいう。)の予想額の現価に相当する額と積立金の額を合算した額を控除した
      額の二十年に一回の頻度で発生すると予想される（　　A　　）とするものであること。

      \lawitem{二}
      （　　B　　）を考慮するものであり、かつ、（　　C　　）を考慮するよう努めているものであること。
      ただし、リスク分担型企業年金を実施する場合にあっては、予定利率と実績とが乖離することに伴い発生し
      うる危険を考慮するものであること。

      \lawitem{三}
      信頼できるデータ、情報及び手法を用いるものであること。
    \end{lawarticle}
  \end{lawstructure}
\end{fillblank}

\begin{answerbox}
  \begin{enumerate}[A: ]
    \item 最大額
    \item 価格変動リスク
    \item 負債変動リスク
  \end{enumerate}
\end{answerbox}

\begin{explanationbox}
  算定告示

  \url{https://www.mhlw.go.jp/web/t_doc?dataId=85ab5661&dataType=0&pageNo=1}
\end{explanationbox}

\newpage

\subsection{給付区分・実施事業所ごとの掛金の算定}


\begin{InyouJobunbox}
  \begin{itemize}
    \item 特例的扱い第3-1(2)〜(6)
  \end{itemize}
\end{InyouJobunbox}


\begin{shortanswer}[年金2 2024 問題2(5)]
  「確定給付企業年金の財政計算等に係る特例的扱いについて(平成 20 年 9 月 11 日年発第 0911001 号)」
  における第3(財政計算時の特例)の第1項第5号に記載されている、実施事
  業所が増加する場合の特別掛金算定の取り扱いを簡潔に入力しなさい。(400 字以内)
\end{shortanswer}

\begin{answerbox}
  \begin{itemize}
    \item 実施事業所が増加する場合は、当該増加に係る財政計算の基準日における当該実施事業所の過去勤務債務の額（基準日における当該事業所の数理債務の額から当該事業所の増加に伴い資産管理運用機関または基金が受換した資産を控除した額をいう。）について、当該事業所の特別掛金を算定することができる。
    \item 当該実施事業所と給付の内容が同一の集団と同じ償却方法により算定すること。
    \item ただし、予定償却期間、償却割合は別に設定できること。
    \item 実施事業所が増加したとしても、規則第50条に該当していない場合（財政再計算実施事由に該当しない場合）は、財政計算を行わず、当該実施事業所の過去勤務債務の額に係る特別掛金のみを算定することができること。
  \end{itemize}
\end{answerbox}

\begin{explanationbox}
  特例的扱い第3-1(5)。
\end{explanationbox}

\begin{fillblank}[数理人試験 2021 問題1 設問1 改題]
    \noindent \textbf{◯特例的扱い}
  \begin{lawstructure}
    \begin{lawarticle}{第三　財政計算時の特例}{}
      \lawitem{一}
      特別掛金の算定に係る原則的扱い

      \hangedlawitem{(1)}{3}
      財政計算に用いる資産額は、基準日における純資産額
      （資産の評価に数理的評価を用いている場合には数理上資産額）
      から（　　A　　）の額を控除した額とし、規則第46条第1項の規定により
      過去勤務債務の額を算定する場合の積立金の額には、当該資産額を用いること。
      なお、規則第46条第1項第1号、第2号及び第4号の方法により特別掛金を算定する場合において、
      基準日以降における（　　B　　）を見込むことにより算定することができること。

      \hangedlawitem{(2)(3)　（略）}{3}

      \hangedlawitem{(4)}{3}
      前記（１）から（３）に定めるところによらず、 （１）から（３）の過去勤
務債務の額について、次の\textcircled{1}又は\textcircled{2}の方法により実施事業所ごとに配分した
額に基づいて、実施事業所ごとに特別掛金が算定できること。この場合にお
いて、同一給付区分においては、規則第四十六条第一項各号に規定する償却
方法のうち同一の償却方法を用いることとし、後記（５）により算定される
特別掛金を除き、同項第一号に規定する予定償却期間（以下「予定償却期間」
という。）の完了日又は第三号に規定する一定の割合（以下「償却割合」とい
う。）は同一のものとすること。なお、過去勤務債務の額が零を下回る実施事
業所がある場合には、他の実施事業所の過去勤務債務の額から当該下回る額
を控除すること。

      \hangedlawitem{\textcircled{1}}{4}
      過去勤務債務の額を直前の財政検証、前回の財政計算又は当該財政計算
の基準日における（　　C　　）のいずれかの比により
按分する方法

      \hangedlawitem{\textcircled{2}}{4}
      過去勤務債務の額から直前の財政検証、前回の財政計算又は当該財政計
算の基準日における（　　D　　）（当該財政計算の基準日におけ
る額の場合は財政計算前の額とする。）を控除した額について\textcircled{1}の方法又
は直前の財政検証、前回の財政計算若しくは当該財政計算の基準日におけ
る数理債務の額から（　　D　　）と（　　E　　）の合計額を控
除した額の比により按分する方法のいずれかにより各実施事業所に配分
した額に各実施事業所の当該（　　D　　）を加算した額とする方法
      
      （略）

      \hangedlawitem{(5)}{3}
実施事業所が（　　F　　）する場合は、前記（１）から（４）に定めるところによ
らず、当該（　　F　　）に係る財政計算の基準日における当該実施事業所の過去勤務
債務の額（基準日における当該事業所の数理債務から当該事業所の増加に伴
い資産管理運用機関又は基金（以下「基金等」という。）が（　　G　　）を控
除した額をいう。）について、当該事業所の特別掛金を算定することができる
こと。この場合において、当該実施事業所と給付の内容が同一の集団と同じ
償却方法により算定すること。ただし、予定償却期間、償却割合は別に設定
できること。なお、実施事業所が（　　F　　）したとしても、規則第五十条に該当し
ていない場合は、財政計算を行わず、当該実施事業所の（　　G　　）のみを算定することができること。

      （以下略）
    \end{lawarticle}
  \end{lawstructure}
\end{fillblank}

\begin{answerbox}
  \begin{enumerate}[A: ]
    \item 別途積立金
    \item 加入者の数又は加入者の給与の額の変動
    \item 加入者の数、給与または数理債務の額
    \item 特別掛金収入現価
    \item 特例掛金収入現価
    \item 増加
    \item 受換した資産
    \item 過去勤務債務の額に係る特別掛金
  \end{enumerate}
\end{answerbox}

\newpage

\subsection{給付区分毎に資産を区分している場合の特例}

出題例なし

\newpage

\subsection{承継事業所償却積立金を設けている場合の特例}



\begin{InyouJobunbox}
  \begin{itemize}
    \item 特例的扱い第3-4、第5
  \end{itemize}
\end{InyouJobunbox}

\begin{fillblank}[数理人試験 H21 問題1(5) 改題]
  \noindent \textbf{◯特例的扱い}
  \begin{lawstructure}
    \begin{lawarticle}{第五　承継事業所償却積立金}{}
      \lawitem{一}
      趣旨
\\
    承継事業所償却積立金は、（　　A　　）に伴い基金等が受換する資産額
が増加時における当該事業所の（　　B　　）と（　　C　　）の合計額を上回る場合に、当該事業所の積
立金として積み立てる勘定科目であること。

      \lawitem{二}
      （略）

      \lawitem{三}
      とりくずす方法
      \\
      （　　D　　）の結果、承継事業所償却積立金を有する実施事業所が（　　E　　）を拠
出することとなるときは、第三の四(1)及び(2)により控除した額につき、承継事業所償却
積立金をとりくずすこと。
      （略）
     
    \end{lawarticle}
  \end{lawstructure}
\end{fillblank}

\begin{answerbox}
  \begin{enumerate}[A: ]
    \item 実施事業所の増加
    \item 数理債務の額
    \item リスク対応掛金収入現価
    \item 財政計算
    \item 特別掛金又はリスク対応掛金
  \end{enumerate}
\end{answerbox}

\newpage